\section{AuthorityBench Dataset}
\label{app:authoritybench}

\paragraph{Overview.}
\textsc{AuthorityBench} is a synthetic benchmark for evaluating whether
language models can determine if a requested action is authorized under
priority-ordered policies. Each example contains users with allow/disallow
rules, an explicit priority order, a requested action or tool use, and a binary
\texttt{Yes}/\texttt{No} label. We provide both general action authorization
and named tool-use authorization tasks.

\paragraph{Dataset configurations.}
We release four configurations. \textsc{GeneralAuthorityV1} and
\textsc{ToolAuthorityV1} use deterministic query templates.
\textsc{GeneralAuthorityV2} and \textsc{ToolAuthorityV2} preserve the
corresponding V1 policies and labels but replace the query with an
LLM-generated natural-language request. Each row contains \texttt{text},
\texttt{label}, and \texttt{meta\_data}, where the metadata records the
structured query, conflict status, query rendering, and paraphrase verification
result when applicable.

\paragraph{Structured policy generation.}
Let the priority order be
\(\pi=(u_1,\ldots,u_m)\), where \(u_1\) has the highest authority. A query is
represented as
\[
q=\{(c,x_c):c\in C_q\},
\]
where \(C_q\) is the set of queried categories and \(x_c\) is the requested
value for category \(c\). Each user \(u\) is assigned structured allow/disallow
rules, which induce a deterministic decision
\(d_u(q)\in\{\mathsf{Yes},\mathsf{No}\}\). The gold label is determined by the
highest-priority user:
\[
y=d_{u_1}(q).
\]
An example is marked as a conflict when
\(d_{u_j}(q)\neq d_{u_1}(q)\) for at least one \(j>1\). The model must therefore
apply the stated priority order rather than aggregate or majority-vote over
user decisions.

\begin{tcolorbox}[
    enhanced,
    colback=gray!3,
    colframe=black!98,
    boxrule=0.5pt,
    arc=1mm,
    left=1.5mm,
    right=1.5mm,
    top=1mm,
    bottom=1mm,
    title=\textbf{Example: Priority-Aware Tool Authorization},
    fonttitle=\small,
    fontupper=\small
]
\textbf{Task.}\par
Determine whether the requested use of
\texttt{send\_email} is authorized.

\smallskip
\textbf{Priority.}\par
User A \(>\) User B.

\smallskip
\textbf{User A.}\par
Allows sending \emph{schedule information} to a
\emph{coworker}; disallows sending it to a
\emph{family member}.

\smallskip
\textbf{User B.}\par
Disallows sending \emph{schedule information} to a
\emph{coworker}.

\smallskip
\textbf{Query.}\par
Use \texttt{send\_email} to send schedule information
to a coworker.

\smallskip
\textbf{Answer.}\par
\texttt{Yes}. Although User B prohibits the action,
the higher-priority User A explicitly permits it.
\end{tcolorbox}


\paragraph{Generalization-oriented splits.}
The train and test sets differ along controlled generalization axes rather than
forming a random IID split. For general authorization, training uses
\emph{day}, \emph{date}, and \emph{time}, while testing additionally introduces
\emph{month} and \emph{year}. For tool authorization, training uses
\emph{recipient}, \emph{information type}, and \emph{purpose}, while testing
also introduces \emph{day} and \emph{time}.

Tool identities are disjoint across splits. Training uses
\texttt{send\_email}, \texttt{send\_text\_message}, and
\texttt{submit\_online\_form}, whereas testing uses
\texttt{make\_phone\_call}, \texttt{send\_voice\_message}, and
\texttt{start\_live\_chat}. Training examples contain one to three users,
whereas test examples contain up to five users and longer policy descriptions.

\paragraph{Natural-query paraphrasing.}
For V2, the paraphrasing model receives only the structured query conditions
and, for tool-use examples, the tool name. It does not receive the policy,
priority order, or gold label. For example,
\texttt{\{"day":"wednesday","time":"13:00--16:59"\}} may be rendered as
``Perform the action on Wednesday between 13:00 and 16:59.''

We use \texttt{gpt-5.4-mini} for query generation and a separate binary
semantic-verification prompt. A paraphrase is retained only when the verifier
determines that it expresses all and only the original structured conditions.
Paraphrases are cached by structured query so that identical requests can be
reused across different policies.

\begin{table}[t]
\centering
\small
\setlength{\tabcolsep}{3.8pt}
\begin{tabular}{@{}lrrr@{}}
\toprule
Configuration & Train & Test & Total \\
\midrule
\textsc{GeneralAuthorityV1} & 500 & 1,500 & 2,000 \\
\textsc{GeneralAuthorityV2} & 491 & 1,468 & 1,959 \\
\textsc{ToolAuthorityV1}    & 1,000 & 3,000 & 4,000 \\
\textsc{ToolAuthorityV2}    & 991 & 2,893 & 3,884 \\
\bottomrule
\end{tabular}
\caption{Dataset configurations. V2 retains only semantically verified
paraphrases.}
\label{tab:authority-dataset-statistics}
\end{table}

\paragraph{Dataset statistics and evaluation.}
The deterministic V1 configurations contain 6,000 examples, with 2,996
\texttt{Yes} and 3,004 \texttt{No} labels. They include 2,846 conflict and
3,154 non-conflict examples. After semantic filtering, the V2 configurations
contain 5,843 examples.

We use exact-match \texttt{Yes}/\texttt{No} accuracy as the primary metric and
additionally report results by conflict status, user count, category
combination, and tool identity. Because the benchmark is synthetic, it provides
precise labels and controlled generalization tests, but does not capture every
ambiguity present in naturally occurring authorization instructions.
